% Options for packages loaded elsewhere
\PassOptionsToPackage{unicode}{hyperref}
\PassOptionsToPackage{hyphens}{url}
\PassOptionsToPackage{dvipsnames,svgnames,x11names}{xcolor}
%
\documentclass[
  10pt,
  a4paper,
]{article}
\usepackage{amsmath,amssymb}
\usepackage{setspace}
\usepackage{iftex}
\ifPDFTeX
  \usepackage[T1]{fontenc}
  \usepackage[utf8]{inputenc}
  \usepackage{textcomp} % provide euro and other symbols
\else % if luatex or xetex
  \usepackage{unicode-math} % this also loads fontspec
  \defaultfontfeatures{Scale=MatchLowercase}
  \defaultfontfeatures[\rmfamily]{Ligatures=TeX,Scale=1}
\fi
\usepackage{lmodern}
\ifPDFTeX\else
  % xetex/luatex font selection
  \setmainfont[]{Liberation Serif}
  \setsansfont[]{Liberation Sans}
  \setmonofont[]{DejaVu Sans Mono}
\fi
% Use upquote if available, for straight quotes in verbatim environments
\IfFileExists{upquote.sty}{\usepackage{upquote}}{}
\IfFileExists{microtype.sty}{% use microtype if available
  \usepackage[]{microtype}
  \UseMicrotypeSet[protrusion]{basicmath} % disable protrusion for tt fonts
}{}
\makeatletter
\@ifundefined{KOMAClassName}{% if non-KOMA class
  \IfFileExists{parskip.sty}{%
    \usepackage{parskip}
  }{% else
    \setlength{\parindent}{0pt}
    \setlength{\parskip}{6pt plus 2pt minus 1pt}}
}{% if KOMA class
  \KOMAoptions{parskip=half}}
\makeatother
\usepackage{xcolor}
\usepackage[margin=22mm]{geometry}
\setlength{\emergencystretch}{3em} % prevent overfull lines
\providecommand{\tightlist}{%
  \setlength{\itemsep}{0pt}\setlength{\parskip}{0pt}}
\setcounter{secnumdepth}{-\maxdimen} % remove section numbering
\ifLuaTeX
\usepackage[bidi=basic]{babel}
\else
\usepackage[bidi=default]{babel}
\fi
\babelprovide[main,import]{korean}
\ifPDFTeX
\else
\babelfont{rm}[]{Liberation Serif}
\fi
% get rid of language-specific shorthands (see #6817):
\let\LanguageShortHands\languageshorthands
\def\languageshorthands#1{}
\usepackage{xeCJK}
\setCJKmainfont{Noto Serif CJK KR}
\setCJKsansfont{Noto Sans CJK KR}
\setCJKmonofont{Noto Sans Mono CJK KR}
\xeCJKsetup{CJKspace=true}
\usepackage{xurl}
\usepackage{fvextra}
\usepackage{fancyhdr}
\AtBeginDocument{\DefineVerbatimEnvironment{verbatim}{Verbatim}{breaklines=true,breakanywhere=true,fontsize=\small}}
\setlength{\emergencystretch}{4em}
\setlength{\headheight}{15pt}
\pagestyle{fancy}\fancyhf{}
\fancyhead[L]{\small 중성자별 구조와 비방사 진동}
\fancyfoot[C]{\thepage}
\renewcommand{\headrulewidth}{0.3pt}
\setcounter{tocdepth}{2}
\ifLuaTeX
  \usepackage{selnolig}  % disable illegal ligatures
\fi
\usepackage{bookmark}
\IfFileExists{xurl.sty}{\usepackage{xurl}}{} % add URL line breaks if available
\urlstyle{same}
\hypersetup{
  pdftitle={P5 · 누적 음향시간의 독립 확장},
  pdflang={ko-KR},
  colorlinks=true,
  linkcolor={Maroon},
  filecolor={Maroon},
  citecolor={Blue},
  urlcolor={Blue},
  pdfcreator={LaTeX via pandoc}}

\title{P5 · 누적 음향시간의 독립 확장}
\usepackage{etoolbox}
\makeatletter
\providecommand{\subtitle}[1]{% add subtitle to \maketitle
  \apptocmd{\@title}{\par {\large #1 \par}}{}{}
}
\makeatother
\subtitle{과제 큐시트}
\author{}
\date{}

\begin{document}
\maketitle

\setstretch{1.12}
\textbf{선행:} P4　 \textbf{독립 작업:} 3--5시간 (계획치)

\subsection{음향시간을 추가하면서 원래 별을 바꿔야
하는가?}\label{uxc74cuxd5a5uxc2dcuxac04uxc744-uxcd94uxac00uxd558uxba74uxc11c-uxc6d0uxb798-uxbcc4uxc744-uxbc14uxafd4uxc57c-uxd558uxb294uxac00}

첫 저장 반경이 0이 아니다. 누적 적분을 그곳에서 0으로 시작하면 무엇을
생략하는지 먼저 그림이나 문장으로 설명하자.

먼저 한두 문장으로 답한 뒤 아래 설명을 읽자. 첫 예측이 틀렸거나 아직
모르겠다고 적는 것은 감점하지 않는다.

\subsection{필요한 관계와 읽을
위치}\label{uxd544uxc694uxd55c-uxad00uxacc4uxc640-uxc77duxc744-uxc704uxce58}

\[
\tau(r)=\int_0^r\frac{e^{\lambda(r')-\nu(r')}}{c\sqrt{\chi(r')}}\,dr'.
\] \(0<r<r_0\)에서는 첫 integrand를 유지하고, 이후에는 저장된 radial
점의 trapezoid rule을 적용한다.

\begin{itemize}
\tightlist
\item
  \path|extensions/P5/cumulative_acoustic_time.py|
\item
  \path|notes/chapters/ch12.md|
\end{itemize}

\subsection{계산으로
옮기기}\label{uxacc4uxc0b0uxc73cuxb85c-uxc62euxae30uxae30}

고급 solver를 바꾸지 않는 작은 진단 확장이다. 독립 폴더의 P5 함수만
작성하고 기존 배경과 진동식은 그대로 둔다.

\begin{enumerate}
\def\labelenumi{\arabic{enumi}.}
\tightlist
\item
  현재 독립 export 안의 P5 slot만 수정한다.
\item
  중심 누락구간의 hold와 이후 trapezoid 규칙을 구현한다.
\item
  상수 integrand와 비균일 격자의 비상수 integrand를 각각 확인한다.
\item
  잘못된 chi/반경 입력을 거절하고 solver source hash를 보존한다.
\end{enumerate}

\newpage

\subsection{실행해서 무엇을 볼
것인가}\label{uxc2e4uxd589uxd574uxc11c-uxbb34uxc5c7uxc744-uxbcfc-uxac83uxc778uxac00}

패키지 루트에서 실행한다. \texttt{-\/-plan}은 경로만 확인하며 계산
성공을 뜻하지 않는다.

\begin{verbatim}
python tools/run_assignment.py P5 --plan
python tools/run_assignment.py P5 --run \
  --workspace learner --output ../P5_result.json
\end{verbatim}

\subsection{결과를
설명하기}\label{uxacb0uxacfcuxb97c-uxc124uxba85uxd558uxae30}

배경을 읽기만 하는 진단이다. 중심 누락 구간의 규칙과 각 반경 간격을
명시하고, 상수 integrand 하나로 일반 적분 구현을 검증했다고 말하지
않는다.

\textbf{남길 증거.} cumulative τ CSV와 metadata; 두 독립 적분 기준의
시험 결과; 허용 변경 경로와 이전 solver 불변 receipt

\textbf{판단 기준.} 2점 예: 비상수 g의 비균일 격자에서도 중심 hold와
trapezoid를 독립 계산으로 확인한다. 최초 예측의 정오보다, 결과를 본 뒤
이유를 고치고 다른 입력에도 적용하는지를 본다.

\subsection{막혔을 때}\label{uxb9c9uxd614uxc744-uxb54c}

\textbf{힌트 1.} 처음 저장된 반경이 0이 아니므로 초기 τ를 0으로 두면 안
된다.

\textbf{힌트 2.} 간격이 균일하다고 가정하지 말고 각 Δr을 사용한다.

\subsection{다음 만남에 다시
묻기}\label{uxb2e4uxc74c-uxb9ccuxb0a8uxc5d0-uxb2e4uxc2dc-uxbb3buxae30}

다음 회차에는 고정 profile의 lapse에 상수를 더한 경우 음향시간이 어떻게
바뀌는지 설명한다. 새 고유모드 계산과는 구별한다. 먼저 짧게 답한 뒤
본문과 자신의 기록을 열어 확인한다.

\textbf{범위.} 선형 sample 보간의 구적 규칙이며 실제 음파 mode
spectrum을 새로 계산하는 solver가 아니다.

\end{document}
